% 由 work/tools/gen_glossary.py 自动生成，勿手改
\chapter*{附录 A\quad 术语对照表}
\addcontentsline{toc}{chapter}{附录 A\quad 术语对照表}

本表给出本书日文底本用语—英文原著用语—中译用语的对应，按英文字母序排列。
译名以此为定准；正文中首次出现的重要术语，另在脚注或批注中说明。
括号内为原书特殊用法或易混之处。

\begin{longtable}{@{}>{\raggedright\arraybackslash}p{0.27\textwidth}>{\raggedright\arraybackslash}p{0.22\textwidth}>{\raggedright\arraybackslash}p{0.17\textwidth}>{\raggedright\arraybackslash}p{0.22\textwidth}@{}}
\toprule
\textbf{英文} & \textbf{日文底本} & \textbf{中译} & \textbf{备注} \\
\midrule
\endhead

1-complement & 1-補足 & 1-补足 & p.274 注意 7.35 [Sho92] \\
1-cycle & 1-サイクル & 1-循环 & 原书「1-サイクル」；与「0-サイクル」同族，保留数字+循环 \\
2-torsion point & 2分点 & 2-挠点 & 原书「2分点」即二阶挠点 \\
a cyclic group acting freely in codimension 1 & 余次元1で自由に作用する巡回群 & 在余维 1 处自由作用的循环群 & ch05b:44/45、ch05c:40「自由循环群」词序误读为 free cyclic group（报告 A3） \\
absolute version & 絶対版 & 绝对版 &  \\
abundance & アバンダンス & 丰度 & 7.74（glossary 丰度猜想） \\
abundance conjecture & アバンダンス予想 & 丰度猜想 & 另译「丰富猜想」，全书统一用丰度猜想 \\
abuse of notation & 記号の濫用 & 记号滥用 & 定理 4.39 証明；M\_r 兼作線束 \\
accumulate & 集積 & 聚积 & 极射线在半空间中聚积 \\
adjunction formula & 随伴公式 & 伴随公式 & 随伴→伴随 \\
affine coordinate system & アフィン座標系 & 仿射坐标系 &  \\
algebraic space & 代数空間 & 代数空间 & 4.66 \\
algebraic valuation & 代数的付値 & 代数赋值 & ch02b 注意 2.23 \\
algebraic variety & 代数多様体 & 代数簇 & 多様体在代数几何语境译“簇” \\
ample & 豊富 & 丰富 &  \\
analytic branch & 解析的分枝 & 解析分支 & 定理 4.15 \\
analytic coordinates & 解析座標 & 解析坐标 & 定理 4.57 \\
analytic family & 解析族 & 解析族 & 定理 4.39 証明；P\_i(t) の解析族 \\
analytic space & 解析多様体 & 解析空间 & 本書语境为正规复解析空间 \\
anharmonic ratio; cross ratio & 非調和比 & 非调和比（交比） & 注意 4.9；正文首见用“交比” \\
anti-canonical & 反標準 & 反标准 &  \\
anti-canonical bundle & 反標準束 & 反标准丛 &  \\
arbitrary & 任意の & 任意的 &  \\
arithmetic genus & 算術種数 & 算术亏格 & 定理 4.7 証明；英文作 arithmetic genus \\
Artin algebra & Artin 多元環 & Artin 代数 & 记号 4.52 \\
assume & と仮定する & 假设 & 动词；名词「仮定」同译「假设」，本书不按词性另造译名 \\
assumption & 仮定 & 假设 & 名词；动词 assume＝仮定する，中文同为「假设」 \\
at most & 高々 & 至多 & 定理 3.47 \\
automorphism & 自己同型 & 自同构 &  \\
band region & 帯状領域 & 带状区域 & 0<aq-rp<1 的带状区域 \\
base change & 底変換 & 基变换 & 不用「底变换」 \\
base locus & 固定点集合 & 基点集 & 原書作「固定点集合」；英文 base locus；亦作「基点集合」 \\
base-point-free linear system & 自由線形系 & 无基点线性系 & ch02c 推论 2.33 \\
basepoint-free & 基点自由 & 无基点 &  \\
basepoint-free theorem & 基点自由性定理 & 无基点定理 &  \\
be nothing but & 他ならない & 不外是 & 「…に他ならない」 \\
bend-and-break lemma & 折曲補題 & 折弯引理 & 原书作「折曲補題」；BRIEF 已定为「折弯引理」，中文依 Mori 的 bend-and-break 通译 \\
big & 巨大 & 大 & 术语作「大」；nef and big＝nef 且大；不用「巨大」 \\
birational geometry & 双有理幾何学 & 双有理几何 & 书名 \\
birational geometry of pairs & 対の双有理幾何 & 对的双有理几何 & 3.30 标题 \\
birational invariant & 双有理不変量 & 双有理不变量 &  \\
birational map & 双有理写像 & 双有理映射 &  \\
birational transformation & 双有理変換 & 双有理变换 & 亦作 birational transform \\
birationally equivalent & 双有理同値 & 双有理等价 &  \\
biregular & 双正則 & 双正则 &  \\
blow-down & ブローダウン & 爆破收缩 &  \\
blow-up & ブローアップ & 爆破 &  \\
branch & 枝 & 枝 & 定義 4.6 \\
branch / ramification & 分岐 & 分歧 & 与「分岐点（fork＝分叉点）」区分；分岐被覆＝分歧覆盖 \\
branched cover & 分岐被覆 & 分歧覆盖 & 定义 2.50、2.52 \\
branched cyclic cover & 分岐巡回被覆 & 分歧循环覆盖 & ch02d 定义 2.50/2.52 \\
bundle & 束 & 丛 & 標準束→标准丛；注意与“層(层)”区分 \\
by the way & ところで & 顺带一提 &  \\
Calabi–Yau variety & カラビ・ヤウ多様体 & 卡拉比–丘簇 &  \\
canonical bundle & 標準束 & 标准丛 &  \\
canonical class & 標準類 & 标准类 &  \\
canonical divisor & 標準因子 & 标准除子 &  \\
canonical flop & 標準フロップ & 标准翻模 &  \\
canonical model & 標準モデル & 标准模型 &  \\
canonical quotient singularity & 標準商特異点 & 标准商奇点 & 定理 7.19 \\
canonical ring & 標準環 & 标准环 &  \\
canonical singularity & 標準特異点 & 标准奇点 &  \\
Cartier divisor & カルティエ因子 & Cartier 除子 & Cartier/Weil 保留拉丁字母，不音译 \\
ch02a:205/211/217、ch02b:6、ch02c:97/113/115/177/194/228、ch02d:97/113/115/177/194、ch02e:87/89/91/96/107/109 & 第 N 步 & 步骤 N & 页图「段階」；glossary 段階=步骤；ch03/ch06 用「步骤 N」 \\
chain & 鎖 & 链 & 定義 4.6；长度＝顶点数 \\
characteristic & 標数 & 特征 & 特征 0 \\
claim & 主張 & 断言 & 证明内以「主張 3.20」形式出现，不用「主张」以免与日常义混淆；亦作 assertion \\
class of pairs & 対のクラス & 对之类 & 3.29、3.31 \\
closed convex cone & 閉凸錐 & 闭凸锥 &  \\
closed point & 閉点 & 闭点 &  \\
closed subscheme & 閉部分スキーム & 闭子概形 & 原书「閉部分スキーム」，依 scheme=概形 \\
closed subvariety & 閉部分多様体 & 闭子簇 &  \\
closure & 閉包 & 闭包 & NE 的闭包 \\
coarse moduli space & 粗モジュライ空間 & 粗模空间 & 7.71 \\
codimension & 余次元 & 余维数 & 引理 4.53 \\
coefficient & 係数 & 系数 &  \\
coherent sheaf & 連接層 & 凝聚层 &  \\
cohomology & コホモロジー & 上同调 &  \\
cokernel & 余核 & 余核 &  \\
commutative diagram & 可換図式 & 交换图式 & 3.31 步骤 1 \\
compact & コンパクト & 紧 &  \\
complete & 完備 & 完备 &  \\
complete linear system & 完備線形系 & 完备线性系 &  \\
completion & 完備化 & 完备化 &  \\
complex & 複体 & 复形 &  \\
complex analytic space & 複素解析空間 & 复解析空间 & 定义 4.59 \\
compound Du Val & 複合 Du Val & 复合 Du Val / cDV & ch05b:229 作「复合 Du Val」、ch05c 全用「cDV」，建议首见全称后统一 cDV（报告 B4） \\
conclusion & 結論 & 结论 &  \\
condition & 条件 & 条件 &  \\
cone & 錐 & 锥 &  \\
cone of curves & 曲線の錐 & 曲线锥 & 記号 NE(X) \\
cone theorem & 錐定理 & 锥定理 &  \\
conic bundle & 二次曲線束 & 二次曲线丛 & bundle＝丛，不用「二次曲线束」 \\
conjecture & 予想 & 猜想 & 環境 conjecture \\
connected & 連結 & 连通 &  \\
contraction & 収縮 & 收缩 &  \\
contraction morphism & 収縮射 & 收缩态射 & 正文 10 处；morphism＝态射，故作「收缩态射」 \\
contraction of an extremal cone face & 端錐面の収縮 & 极锥面的收缩 & 原書 p.32 末 \\
convergent power series & 収束べき級数 & 收敛幂级数 & 定理 4.61 以下 \\
convex cone & 凸錐 & 凸锥 &  \\
corollary & 系 & 推论 & 環境 corollary \\
counterexample & 反例 & 反例 & 環境 counterexample \\
covering & 被覆 & 覆盖 &  \\
cusp & カスプ & 尖点 & ノート 4.8；与 cuspidal cubic 的“尖点”同一词，靠上下文区分 \\
cuspidal cubic & 尖点付三次曲線 & 尖点三次曲线 & ノート 4.8 \\
cycle & サイクル & 圈 & 定義 4.6；n＝2 为两顶点两平行边 \\
cyclic & 巡回 & 循环 &  \\
cyclic cover & 巡回被覆 & 循环覆盖 & §2.4 第 3 步；「不分岐/分岐」作不分歧/分歧 \\
cyclic quotient singularity & 巡回商特異点 & 巡回商奇点 & 定理 4.7(5) \\
decomposition into connected components & 連結成分への分解 & 到连通分支的分解 & 记号 4.52 \\
definition & 定義 & 定义 & 環境 definition；末尾□ \\
definition of intersection number & 交点数の定義 & 交点数的定义 & 原书 1.34 的标题；四种手法并列 \\
deformation & 変形 & 形变 & 変形理論→形变理论 \\
deformation space & 変形空間 & 形变空间 &  \\
deformation space with base condition & 底条件付き変形空間 & 附带底条件的形变空间 & 定理 1.10 步骤 2；原书「底条件付き」，指固定 f(0) 后仍可形变 \\
Del Pezzo fibration & Del Pezzo ファイバー空間 & Del Pezzo 纤维空间 & 定理 1.32 的 D 型；ファイバー=纤维 \\
Del Pezzo surface & Del Pezzo 曲面 & Del Pezzo 曲面 & 保留 \\
desingularization & 脱特異点 & 奇点消解 & 与 resolution of singularities 同义；正文统一「奇点消解」 \\
diagonal & 対角線 & 对角线 & 構成 4.34；C\^\{\}n×C の第 i 対角線 \\
direct image & 順像 & 顺像 & 层论 direct image（顺像层）；与 pushforward 同译「顺像」 \\
discrepancy & 食い違い係数 & 差异系数 & 第3章后半大量使用 \\
discrete valuation ring & 離散付値環 & 离散赋值环 & p.285 \\
divisible group & 可除群 & 可除群 & p.269 μ\_∞ 的性质 \\
divisor & 因子 & 除子 & factor=因子；divisor=除子 \\
divisor / prime divisor & 因子（误作除子） & 除子 / 素除子 & ch02c:208/231/234/258/261/266/267/279、ch07c:197/221；factor 义的「直和因子/公共因子/零因子」不动 \\
divisor bundle & 因子束 & 除子丛 &  \\
divisor over X & 空の因子 & 上空除子 & ch02b 定义 2.24 \\
divisorial contraction & 因子収縮 & 除子收缩 &  \\
divisorial log terminal & 因子ログ端末的 & 除子对数终端 & dlt 类；ログ端末→对数终端已冻结 \\
dlt & dlt & 除子对数终端 & 定義 2.37；全称作「除子对数终端」（divisional log terminal）；原書「因子ログ端末的」按 divisor＝除子 改 \\
dominate & 支配する & 支配 & f 支配爆破 B\_x X \\
Du Val singularity & デュヴァル特異点 & Du Val 奇点 & 也译“杜瓦尔奇点”；正文用 Du Val \\
dual cone & 双対錐 & 对偶锥 &  \\
dual graph & 双対グラフ & 对偶图 &  \\
duality theory & 双対理論 & 对偶理论 &  \\
dualizing sheaf & 双対化層 & 对偶化层 & 正文 11 处；第 5–7 章关键技术词 \\
edge & 辺 & 边 & 定義 4.6 \\
effective & 有効 & 有效 &  \\
eigensheaf & 固有層 & 特征层 & Z/mZ 作用下的特征子层；与「固有」（proper）区别，proper 仍作「固有」；(与 proper=固有 区分) \\
eigenspace & 固有空間 & 特征空间 & 同上；特征层之纤维 \\
eigenvector & 固有ベクトル & 特征向量 & 与「固有」（proper）区分；不用「固有向量」 \\
elimination theory & 消去理論 & 消去理论 & 题辞中「消去理論の雰囲気」；代数学の消去理论 \\
elliptic & 楕円的 & 椭圆的 & 7.79 \\
elliptic singularity & 楕円特異点 & 椭圆奇点 &  \\
elliptic surface & 楕円曲面 & 椭圆曲面 &  \\
embedding dimension & 埋め込み次元 & 嵌入维数 & 定理 4.22 \\
end of proof & 証明終わり & 证毕 & ■ \\
endomorphism & 自己準同型 & 自同态 & 原书「n 倍自己準同型」；同态=準同型，自同态=自己準同型 \\
étale & エタール & 平展 &  \\
étale equivalent & エタール同値 & étale 等价 & ch04e 定義 4.67；勿写「同值」 \\
étale fundamental group & エタール基本群 & 平展基本群 &  \\
Euclidean topology & Euclid 位相 & Euclid 拓扑 & 与 étale 拓扑并列 \\
Euler characteristic & Euler 標数 & Euler 标数 & ch02d 命题 2.57 \\
exact sequence & 完全列 & 正合列 &  \\
example & 例 & 例 & 環境 example；末尾□ \\
exceptional curve & 例外曲線 & 例外曲线 &  \\
exceptional divisor & 例外因子 & 例外除子 & 按 Lead 更正 divisor＝除子（「因子」只用于 factor）；旧记法「例外因子」不再用；亦作 exceptional divisor / curve \\
exceptional set & 例外集合 & 例外集 &  \\
extended dual graph & 拡張双対グラフ & 扩张对偶图 & 定義 4.6；加入非固有曲线 D\_i 作无标号顶点 \\
extremal & 端的 & 极端的 & 定義 3.34 \\
extremal cone & 端錐 & 极锥 &  \\
extremal contraction & 端射線収縮 & 极收缩 & 日文「端射線収縮」；端射線＝极射线 \\
extremal face & 端錐面 & 极锥面 & 英文 extremal face；原书「錐面」指凸锥的面、「端錐面」指由负极射线张成的极面（extremal cone face）；本书统一作极锥面，配套 端射線＝极射线 \\
extremal ray & 端射線 & 极射线 & extremal ray；与“终端奇点(terminal)”区分 \\
extreme point & 端点 & 极点 & 凸几何用语 \\
f-free & f-自由 & f-自由 & 系列 4.14 証明；“O\_Y(-E) は f-自由” \\
f-minimal & f-極小 & f-极小 & 定義 3.48 \\
f-nef & f-ネフ & f-nef & 系列 4.14 証明；保留拉丁缩写 \\
face of a cone & 錐面 & 锥面 & 同上，按原书字面译 \\
faithfully flat & 忠実平坦 & 忠实平坦 &  \\
family & 族 & 族 &  \\
family of curves & 曲線族 & 曲线族 &  \\
Fano variety & Fano 多様体 & Fano 簇 & 保留 Fano 不译 \\
fiber & ファイバー & 纤维 &  \\
fiber product & ファイバー積 & 纤维积 & 引理 4.72 \\
fiber space & ファイバー空間 & 纤维空间 & 2.14 (2)「極小モデルのファイバー空間」 \\
fibre contraction & ファイバー収縮 & 纤维收缩 & Fano 收缩情形 \\
finitely generated & 有限生成 & 有限生成 & 原书「有限生成」；修饰环、域等 \\
fixed component & 固定成分 & 固定分支 &  \\
fixed part & 固定部 & 固定部分 & 与「可動部＝移动部分」相对 \\
flat & 平坦 & 平坦 &  \\
flat family & 平坦族 & 平坦族 &  \\
flatness criterion & 平坦性の判定法 & 平坦性判定法 & 定理 4.39 証明；[Har77,III.9.9] \\
flex & 変曲点 & 拐点 & 例 4.36–4.38；三次曲线的拐点 \\
flip & フリップ & 翻转 &  \\
flipped curve & 後フリップ曲線 & 后翻转曲线 & 定義 3.33 \\
flipping contraction & フリップ収縮 & 翻转收缩 & 定義 3.33；亦作 flip contraction \\
flipping curve & 前フリップ曲線 & 前翻转曲线 & 定義 3.33 \\
flop & フロップ & 翻模 & 首次出现加译注：flop，与小收缩互为对偶 \\
fork & 分岐点 & 分叉点 & 定義 4.6；对偶图的树中连接至少 3 条边的顶点；与「分岐（ramification，分歧）」无关，故不用「分歧点」 \\
fundamental group & 基本群 & 基本群 &  \\
g-trivial & g-自明 & g-平凡 & 定理 7.11 的 T \\
Galois & ガロワ & 伽罗瓦 &  \\
general fibre & 一般ファイバー & 一般纤维 & 正文 13 处 \\
general intersection theory & 一般交叉理論 & 一般相交理论 & 1.34 第三手法；交差＝相交 \\
general member & 一般元 & 一般元 & member＝元；general member 一般元，与 general point（一般点／泛点）区分 \\
general point & 一般点 & 一般点 & 引理 4.50 证明；页图作「生成点」，按英文本改译 \\
general position & 一般の位置 & 一般位置 &  \\
general section & 一般切断 & 一般截面 & ch04d 定理 4.57 证明段階 1 \\
generic fiber & 生成ファイバー & 一般纤维 & 原书与「一般ファイバー」并用，统一作一般纤维 \\
generic flatness & 生成的平坦性 & 生成平坦性 & [Mat69, Thm. 52]；原书「生成的平坦性」 \\
generic point & 生成点 & 泛点 & 与「一般点」区分，定義 3.38(4)(5)；(原书「生成点」；中文通行作泛点) \\
genus & 種数 & 亏格 & 正文首次出现加注「亏格（genus）」 \\
geometric invariant theory (GIT) & 幾何学的不変量理論 & 几何不变量论 & 7.72 \\
geometric point & 幾何学的点 & 几何点 & 原书「幾何学的点」；概形论用语 \\
global ramified covering & 大域分歧覆盖 & 大域分歧覆盖 & p.275 注意 7.37 \\
global section & 大域切断 & 全局截面 & H\^\{\}0 \\
globally analytically Q-factorial & 大域解析的 Q-分解的 & 全局解析 Q-分解的 & ch02b 例 2.17 \\
globally generated & 大域的に生成される & 全局生成的 &  \\
graded & 次数付 & 分次 & 分次代数 \\
half-space & 半空間 & 半空间 & K<0 给出的半空间 \\
hence & 故に & 故 & 比“従って”更正式 \\
higher direct image & 高次順像 & 高次顺像 & ch02d 定理 2.48 \\
hold & が成り立つ & 成立 &  \\
homogeneous coordinates & 同次座標 & 齐次坐标 & 原书「同次座標」 \\
homogeneous part & 斉次部分 & 齐次部分 & 4.24(3)、4.25 \\
homology & ホモロジー & 同调 & 正文 33 处 \\
homomorphism & 準同型 & 同态 &  \\
horizontal & 水平 & 水平的 & 定义 7.1 后，B 的水平性 \\
hyperbolic & 双曲的 & 双曲的 &  \\
hyperplane section & 超平面切断 & 超平面截面 & 例 7.61；切断→截面 \\
hypersurface & 超曲面 & 超曲面 &  \\
hypersurface singularity & 超曲面特異点 & 超曲面奇点 &  \\
ideal sheaf defining & 定義イデアル層 & 定义理想层 & 命题 1.41；定義=定义 \\
image & 像 & 像 &  \\
indecomposable & 直積でない & 不可分解的 & 7.75 \\
indeterminacy point & 不確定点 & 不定点 & 原书「f\^\{\}\{-1\} の不定点」；正文用「不定点」，与「不定点解消」一致 \\
injective & 単射 & 单射 & 正文 19 处 \\
integer coefficients & 整係数 & 整系数 & 原书「整係数」 \\
integrally closed & 整閉 & 整闭 &  \\
intersect transversally & 横断的に交叉する & 横截相交 & 定理 1.38 证明；横断的=横截 \\
intersection & 交差 & 相交 &  \\
intersection cohomology & 交叉コホモロジー & 交叉上同调 &  \\
intersection form & 交差形式 & 相交形式 &  \\
intersection matrix & 交差行列 & 相交矩阵 &  \\
intersection number & 交点数 & 相交数 &  \\
intersection point & 交点 & 相交点 &  \\
intersection product & 交差積 & 相交积 &  \\
inversion of adjunction & 逆同伴 & 伴随反转 & 日文“逆同伴”字面为逆伴随，中文取“伴随反转” \\
irreducible & 既約 & 不可约 &  \\
irreducible component & 既約成分 & 不可约分支 &  \\
isomorphism & 同型 & 同构 &  \\
kernel & 核 & 核 &  \\
Kodaira dimension & 小平次元 & 小平维数 & 記号 κ \\
Kodaira vanishing theorem & 小平消滅定理 & 小平消没定理 & 消滅→消没 \\
lattice bundle & 格子束 & 格丛 & 例 1.8 中 V=M⊕M 的商 V/L；原书「格子束」，依 bundle=丛 \\
lattice point & 格子点 & 格点 & (x,y)-平面的格点 \\
lemma & 補題 & 引理 & 環境 lemma \\
Leray spectral sequence & Leray スペクトル系列 & Leray 谱序列 & ch02e 段階 4 \\
line bundle & 直線束 & 线丛 &  \\
linear system & 線形系 & 线性系 &  \\
linearly equivalent & 線形同値 & 线性等价 & 記号 ≡ 或 ∼ \\
local Nullstellensatz & 局所零点定理 & 局部零点定理 & 定義 4.29 \\
local ring & 局所環 & 局部环 &  \\
local system & 局所系 & 局部系统 & 绕除子一周单值为 ξ\^\{\}r 的局部系统 \\
locally analytically Q-factorial & 局所解析的 Q-分解的 & 局部解析 Q-分解的 & ch02b 例 2.17 \\
locally free sheaf & 局所自由層 & 局部自由层 &  \\
log canonical & ログ標準 & 对数标准 &  \\
log category & ログ圏 & 对数范畴 & 与「ログ因子=对数除子」同族；指对数几何的工作范畴 \\
log Del Pezzo surface & ログ Del Pezzo 曲面 & 对数 Del Pezzo 曲面 & 定理 3.47(3) \\
log divisor & ログ因子 & 对数除子 &  \\
log Kodaira dimension & 対数的小平次元 & 对数小平维数 &  \\
log minimal model & ログ極小モデル & 对数极小模型 & 定理 3.47(1) \\
log resolution & 対数特異点解消 & 对数奇点解消 & p.275 命题 7.34 证明 \\
log ruled surface & ログ線織曲面 & 对数直纹曲面 & 定理 3.47(2) \\
log singularity & ログ特異点 & 对数奇点 &  \\
log smooth & ログ平滑 & 对数光滑 &  \\
log terminal & ログ端末 & 对数终端 &  \\
log terminal (lt) & ログ端末的 & 对数终端 & ch02c 定义 2.34 \\
logarithmic & 対数的 & 对数的 & ログ的正式写法 \\
map & 写像 & 映射 &  \\
maximal ideal & 極大イデアル & 极大理想 & 原书「極大イデアル」 \\
minimal model & 極小モデル & 极小模型 &  \\
minimal model program (MMP) & 極小モデルプログラム & 极小模型纲领 &  \\
minimal resolution & 極小特異点解消 & 极小奇点消解 &  \\
minimal ruled surface & 極小線織曲面 & 极小直纹曲面 & 原书「極小線織曲面」 \\
miniversal deformation & 半普遍変形 & 半普遍形变 & 定义 4.59 及脚注 *1 \\
moduli problem & モジュライ問題 & 模空间问题 & §2.3 標準的特異点条目 \\
moduli space & モジュライ空間 & 模空间 & p.284 7.56 \\
monodromy & モノドロミー & 单值 &  \\
monomial & 単項式 & 单项式 & ch04d 定義 4.56 \\
Mori's program & 森プログラム & 森纲领 &  \\
morphism & 射 & 态射 &  \\
movable divisor & 移動因子 & 移动除子 &  \\
moving part & 可動部 & 移动部分 & f\^\{\}*|mD| = |L|（移动部分）+ 固定部分 \\
multiplicity & 重複度 & 重数 &  \\
Nakai--Moishezon criterion & 中井--Moishezon 判定法 & 中井–Moishezon 判别法 & 定理 1.37；原书「中井-Moishezon 判定法」，判定法=判别法 \\
Nakayama lemma & 中山の補題 & 中山引理 & 亦作 Nakayama's lemma；引理=lemma；引理 4.55 证明、命题 1.41 证明 \\
namely & 即ち & 即 &  \\
nef & ネフ & nef & 保留拉丁，不译作「数值有效」 \\
nef exceptional divisor lemma & ネフ例外因子補題 & nef 例外除子引理 & 引理 3.39 \\
negative section & 負断面 & 负截面 & 例 1.46；断面=截面 \\
neighbourhood & 近傍 & 邻域 & 正文 35 处 \\
nodal cubic & 結節点付三次曲線 & 结点三次曲线 & ノート 4.8 \\
non-vanishing theorem & 非消滅定理 & 非消没定理 &  \\
nonsingular & 非特異 & 非奇异 &  \\
nonsingular variety & 非特異多様体 & 非奇异簇 &  \\
normal & 正規 & 正规 &  \\
normal analytic variety/space & 正規解析多様体 & 正规解析流形 & 原書 p.133；對應「複素多様体＝复流形」；ch03d/03e 現作「正规解析空间」，待術語表定案 \\
normal bundle & 法束 & 法丛 & 定理 1.32 E5；束=丛 \\
normal crossing & 正規交叉 & 正规交叉 & 正规交叉除子 \\
normal projective scheme & 正規射影的スキーム & 正规射影概形 & 定義 3.33（英文作 normal scheme，从日文） \\
normalization & 正規化 & 正规化 &  \\
notation & 記号 & 记号 & 環境 notation \\
note that & なお & 另外 &  \\
numerical & 数値的 & 数值的 &  \\
numerical dimension & 数値的次元 & 数值维数 &  \\
numerical equivalence & 数値的同値 & 数值等价 & ≡；命题 2.57 的结论「只依赖数值等价类」 \\
numerically dlt & 数値的dlt & 数值 dlt & 記法 4.1，dlt 保留英文 \\
numerically f-trivial & 数値的f-自明 & 数值 f-平凡 & 定理 4.5 証明 \\
numerically lc & 数値的lc & 数值 lc & 記法 4.1，lc 保留英文 \\
numerically trivial & 数値的トリヴィアル & 数值平凡 &  \\
obstruction & 障害 & 障碍 &  \\
of finite type & 有限型 & 有限型 &  \\
of general type & 一般型 & 一般型 &  \\
one-parameter family & 1-パラメータ族 & 1-参数族 & p.287 \\
open embedding & 開埋め込み & 开嵌入 & 命题 4.70 \\
open variety & 開多様体 & 开簇 & §2.3 川又端末条目 \\
ordinary blow up & 標準的爆発 & 标准爆破 & 定理 4.57（k=2 与 B\_0X 比较） \\
ordinary double point & 通常 2 重点 & 通常 2 重点 & 定理 1.32 E3；原书「通常 2 重点」，保留“2 重点”说法 \\
ordinary node & 通常結節点 & 通常结点 & 定義 7.57；ch07a 已用 \\
over a field & 体上 & 域上 & ch03e:29「X 为体上正规概形」 \\
pair & 対 & 对 & ch03 起大量出现（klt 对、lc 对、对数标准对） \\
pencil & ペンシル & 束 & 原书「ペンシル」；与 bundle=丛 区分（例 1.23(4)「三次曲線の束」） \\
plurigenus & 多重種数 & 多重亏格 & 7.71 \\
pointed & 点付 & 带点 & 命题 4.69「点付代数的スキーム」 \\
pointed scheme & 点付スキーム & 点付概型 & ch04e 命题 4.69 证明 \\
preimage / inverse image & 逆像 & 逆像 &  \\
prime divisor & 素因子 & 素除子 &  \\
principal polarization & 主極化 & 主极化 & 7.75 \\
problem & 問題 & 问题 & 環境 problem \\
projection formula & 射影公式 & 投影公式 & G\_\{m+r\} ≅ G\_m ⊗ L \\
projective morphism & 射影射 & 射影态射 & ch03e:29（2 处）、ch03e:80（2 处） \\
projective variety & 射影多様体 & 射影簇 &  \\
projectively normal & 射影的正規 & 射影正规 & 例 7.61 \\
proof & 証明 & 证明 & 環境 proof；末尾■ \\
proper & 固有 & 固有 &  \\
proper morphism & 固有射 & 固有态射 &  \\
proposition & 命題 & 命题 & 環境 proposition \\
pseudo canonical model & 擬標準モデル & 伪标准模型 & (4\^\{\}c) \\
pseudo minimal model & 擬極小モデル & 伪极小模型 & (5\^\{\}m)，dlt 时即极小模型 \\
pseudo-effective & 擬有効 & 伪有效 &  \\
pullback & 引き戻し & 拉回 &  \\
purely log terminal & 純ログ端末的 & 纯对数终端 & plt 的全称 \\
pushforward & 押し出し & 顺像 & 与 direct image 同义；层论中「顺像层」；不用「推出」（正文「推出」均作动词 derive） \\
put / set & とおく & 记 & 数学写作「记 x=…」 \\
Q-Cartier & Q-カルティエ & Q-Cartier &  \\
Q-factorial & Q-分解的 & Q-分解的 &  \\
Q-factorialization & Q-分解 & Q-分解 & Q-factorial；首次出现加译注说明“Q-分解”即 Q-factorial \\
quadruple point & 4 重点 & 4 重点 & 定理 1.32 E5；与「2 重点」同族 \\
quasi-finite & 擬有限 & 拟有限 & 引理 4.72 \\
quasi-projective & 準射影的 & 拟射影的 &  \\
question & 問 & 问 & 環境 question \\
quotient & 商 & 商 &  \\
quotient by the group acting on the flop & フロップの群による商 & 翻模按群取商 & 原書 p.52 図2.1 図題；figtitle 訳文核対一致 \\
quotient singularity & 商特異点 & 商奇点 &  \\
radical & 根基 & 根 & 记号 4.52「その根基を m とおく」 \\
rank & 階数 & 秩 &  \\
rank one & 階数1の & 秩 1 的 &  \\
rational curve & 有理曲線 & 有理曲线 &  \\
rational singularity & 有理特異点 & 有理奇点 &  \\
rational surface & 有理曲面 & 有理曲面 &  \\
rationality theorem & 有理性定理 & 有理性定理 &  \\
ray type & 射線型 & 射线型 & 定義 3.34 \\
reduced & 被約 & 约化 & reduced divisor＝约化除子、reduced fibre＝约化纤维、reduced scheme＝约化概形 \\
reduced fraction & 既約分数 & 既约分数 & r=u/v 既约 \\
reduced structure & 被約構造 & 约化结构 & 基点集赋以约化结构 \\
reduction & 還元 & 约化 & 主译「约化」；问题化简语境作「归约」（如归约到低维），半稳定约化定理沿用「约化」 \\
reflexive sheaf & 反射的層 & 自反层 & L\^\{\}\{[i]\} 的定义 \\
regular & 正則 & 正则 &  \\
regular function & 正則有理関数 & 正则有理函数 & 4.66 脚注 *3 \\
relative 1-cycle & 相対的 1-サイクル & 相对 1-循环 & ch02b 例 2.16 \\
relative Kleiman criterion & 相対 Kleiman 判定法 & 相对 Kleiman 判定法 & 定理 1.44 \\
relative minimal model & 相対極小モデル & 相对极小模型 &  \\
relative Picard number & 相対的 Picard 数 & 相对 Picard 数 & ch02b 例 2.16 \\
relative version & 相対版 & 相对版 & 与「絶対版＝绝对版」相对 \\
relatively minimal & 相対極小 & 相对极小 & 定義 3.48 \\
relativization & 相対化 & 相对化 &  \\
remark & 注意 & 注 & 環境 remark；末尾□ \\
resolution of singularities & 特異点解消 & 奇点消解 & 与 desingularization 同义；正文统一「奇点消解」 \\
retract & レトラクト & 收缩核 & 引理 4.13；与“収縮（contraction）”区分 \\
rigidity lemma & 剛性補題 & 刚性引理 & 原书小标题；译作刚性引理 \\
round-down & 切り下げ & 取下整 & ⌊mΔ⌋ \\
round-up & 切り上げ & 取上整 & ⌈D⌉ \\
ruled surface & 線織曲面 & 直纹曲面 & 定理 3.47(2) \\
run (an MMP) & 走らせる & 运行 / 施行 & ch05d:22「跑 MMP」为口语，改「运行」（报告 A9） \\
running the MMP & 極小モデルの実行 & 极小模型的运行 & 原書「実行」；亦作「实施」，正文用「运行」 \\
satisfy & を満たす & 满足 &  \\
scheme & スキーム & 概形 & ch03d 定義 3.33 \\
section & 切断 & 截面 & 構成 4.34；大域切断已收作「全局截面」 \\
self-intersection number & 自己交点数 & 自相交数 &  \\
semi-ample & 半豊富 & 半丰富 &  \\
semi-local ring & 半局所環 & 半局部环 & 记号 4.52 \\
semi-log canonical & 半ログ標準的 & 半对数标准的 & slc，注意 7.5 \\
semi-stable & 半安定 & 半稳定 &  \\
semi-stable reduction & 半安定還元 & 半稳定约化 & 7.17/7.19 \\
semi-stable reduction theorem & 半安定還元定理 & 半稳定约化定理 &  \\
semi-universal deformation space & 半普遍変形空間 & 半普遍形变空间 & 章導言 §4.5 \\
sheaf & 層 & 层 & 連接層→凝聚层 \\
simple elliptic & 単純楕円的 & 单纯椭圆型 & ノート 4.8（1）型 \\
simple normal crossing & 単純正規交差 & 简单正规交叉 & 正文 3 处；与 snc 并用（第 2、4 章） \\
simple normal crossing divisor & 単純正規交叉因子 & 单纯正规交叉（snc）除子 & ch02c；建议统一为 snc 除子 \\
simple singularity & 単純特異点 & 单纯奇点 & 注意 4.21(2) \\
simply connected & 単連結 & 单连通 &  \\
simultaneous resolution & 同時特異点解消 & 同时奇点消解 &  \\
singular cohomology & 特異コホモロジー & 奇异上同调 &  \\
singular fiber & 特異ファイバー & 奇异纤维 &  \\
singularity & 特異点 & 奇点 & 特異点解消→奇点消解 \\
small contraction & 小収縮 & 小收缩 &  \\
smoothability & 平滑化可能性 & 光滑化可能性 & 7.78 \\
snake lemma & 蛇補題 & 蛇引理 & 引理 4.55 证明 \\
socle & ソークル & 基座 & 记号 4.52；首次出现保留 (socle) \\
spectral sequence & スペクトル系列 & 谱序列 & 正文 12 处 \\
spectrum & スペクトル & 谱 & Spec V，p.285 \\
splitting & splitting & 分裂 & ch02e 段階 1 \\
stable curve & 安定曲線 & 稳定曲线 & 定義 7.57 \\
stable curve family & 安定曲線族 & 稳定曲线族 & 定義 7.57 \\
stack & スタック & 栈 & p.162 证明中；同时奇点消解的栈 \\
standard terminalization & 基準端末化 & 基准终端化 & 与本書「標準=canonical」区分 \\
Stein factorization & Stein 分解 & Stein 分解 & 保留人名 \\
Step & 段階 & 第 N 步 & 原书「第 0 段階/第 1 段階…」译作「第 0 步/第 1 步」 \\
string theory & 弦理論 & 弦论 & 7.78 \\
subalgebra & 部分多元環 & 子代数 & 命题 4.54 证明 \\
subring & 部分環 & 子环 & 原书「部分環」；「部分スキーム」译闭子概形而「部分環」作子环，统一按现代汉语习惯 \\
subvariety & 部分多様体 & 子簇 &  \\
sufficiently general & 十分一般 & 充分一般 &  \\
summary & 要約 & 概要 & 原书「要約 1.31」指纲领各步骤的概要 \\
support & 台 & 支集 & (凸几何中的 supporting 仍作「支撑」) \\
support function & 支持関数 & 支撑函数 & N 上的支撑函数 \\
support of a divisor & 因子の台 & 除子的支集 & 台→支集 \\
surjective & 全射 & 满射 & 正文 60 处 \\
terminal & 端末 & 终端 & 端末特異点→终端奇点 \\
terminal flop & 端末フロップ & 终端翻模 &  \\
terminal lemma & 端末補題 & 终端补题 & 原书 p.214「端末補題」，ch05c:79 已译「终端补题」，建议收入 glossary（报告 B5） \\
terminal singularity & 端末特異点 & 终端奇点 &  \\
terminalization & 端末化 & 终端化 &  \\
that is & すなわち & 亦即 &  \\
theorem & 定理 & 定理 & 環境 theorem；无证明时末尾■ \\
therefore & 従って & 因此 & 表因果 \\
threefold & 三次元多様体 & 三维簇 & 正文 23 处 \\
thus & かくして & 于是 &  \\
topological method & 位相的手法 & 拓扑手法 & 1.34 第四手法；位相=拓扑 \\
topologically constructible sheaf & 位相的構成可能層 & 拓扑可构造层 & 第 2 步 \\
torsion subgroup & 挠子群 & 挠子群 & p.269 极大挠子群 \\
total discrepancy & 全食い違い係数 & 全差异 & 定義 2.28；inf 取遍 X 上空一切除子 \\
totaldiscrep & totaldiscrep & 全差异系数 & ch05c:143/152 保留英文 totaldiscrep，建议 glossary 注明「与 discrep 相对，允许非例外除子」（报告 §5） \\
trace map & 跡写像 & 迹映射 & 正文 3 处（ch05d 命题 5.77 引入） \\
unique & 一意的 & 唯一的 &  \\
unit & 単位 & 单位 & 4.24(2)；既有表未收 \\
unramified & 不分岐 & 非分歧 &  \\
unramified covering & 非分歧被覆 & 不分歧覆盖 & 原书「不分岐被覆」；分歧=分歧 \\
vanishing & 消滅 & 消没 &  \\
vanishing theorem & 消滅定理 & 消没定理 &  \\
variety & 多様体 & 簇 & 複素多様体→复流形；実多様体→实流形 \\
versal deformation & 準普遍変形 & 准普遍形变 & 定义 4.59；本书 versal = 准普遍 \\
vertex & 頂点 & 顶点 & 定義 4.6 \\
very ample & 非常に豊富 & 非常丰富 &  \\
weak canonical model & 弱標準モデル & 弱标准模型 & lc 且 Δ 有效时 \\
weak pseudo canonical model & 弱擬標準モデル & 弱伪标准模型 & 定義 3.50 \\
weight & 重み & 权 & 定义 4.56；A\^\{\}n 的权 w \\
weighted blow up & 重み付爆発 & 加权爆破 & 定义 4.56 \\
weighted projective space & 重み付き射影空間 & 加权射影空间 &  \\
Weil divisor & ヴェイユ因子 & Weil 除子 & Cartier/Weil 保留拉丁字母，不音译 \\
\bottomrule
\end{longtable}

\begin{flushleft}
\small 共收录 463 条。
\end{flushleft}
